\documentclass[../../../include/open-logic-section]{subfiles}

\begin{document}

\olfileid{sth}{story}{urelements}
\olsection{أندخل العناصر الأولية؟}

سنستنبط في الفصول الآتية مسلّمات من التصور التراكمي-التكراري للمجموعة؛
ولكن ينبغي قبل المضي أن نبين أمر \emph{العناصر الأولية}.

لم تجز الصورة في~\olref[approach]{sec} إنشاء مجموعة جديدة إلا من
\emph{مجموعات} سبقتها. وقد نريد أن نجعل بعض \emph{اللامجموعات}، كالبقر
والخنازير وحبات الرمل وسواها، !!{element} في مجموعات. فنبدأ بعناصر
أساسية نسميها \emph{العناصر الأولية}، ثم نكوّن في كل مرحلة~${\OLId{latin-looped}{S}}$ «كل
مجموعة ممكنة» من عناصر أولية أو من مجموعات نشأت في مراحل سبقت~${\OLId{latin-looped}{S}}$،
مجتمعة أو منفردة. وتقارب الصورة عندئذ الرسم الآتي:
\begin{center}
	\begin{tikzpicture}[scale=0.6]
	\tikzset{cut_here/.style={densely dotted}}
	\draw (-4,6) -- (-1, 0) -- (1,0) -- (4,6); 
	\draw[cut_here] (-5,8)--(-4,6);
	\draw[cut_here] (5,8)--(4,6);
	\draw(-1.5, 1)--(1.5, 1);
	\draw(-2, 2)--(2, 2);
	\draw(-2.5, 3)--(2.5,3);
	\draw(-3, 4)--(3, 4);
	\draw(-3.5, 5)--(3.5, 5);
	\draw(-4, 6)--(4, 6);
	\node[label] at (2, 0) {\small 0};
	\node[label] at (2.5, 1) {\small 1};
	\node[label] at (3, 2) {\small 2};
	\node[label] at (3.5, 3) {\small 3};
	\node[label] at (4, 4) {\small 4};
	\node[label] at (4.5, 5) {\small 5};
	\node[label] at (5, 6) {\small 6};
	\end{tikzpicture}
\end{center}
فلا بد الآن من الفصل في السؤال: \emph{أنسمح بالعناصر الأولية؟}

ومن جهة الفلسفة، إدخال العناصر الأولية معقول، وأقوى دواعيه أن تصبح
نظرية المجموعات \emph{صالحة للتطبيق}. فقد تقرر في
\olref[sfr][siz][]{chap} أن للمجموعتين~${\OLId{latin-looped}{A}}$ و${\OLId{latin-looped}{B}}$ قدرًا واحدًا، أي
$\cardeq{{\OLId{latin-looped}{A}}}{{\OLId{latin-looped}{B}}}$، إذا وفقط إذا قام بينهما تقابل. فإن كان بقر الحقل
وخنازير الحظيرة مجموعتين، أمكن صوغ «عدد البقر كعدد الخنازير» في نظرية
المجموعات. وإن منعنا العناصر الأولية فلم \emph{يكن} البقر والخنازير
مجموعتين، تعذر هذا الصوغ، ولم تنطبق النظرية مباشرة على شيء غير
المجموعات. ولمزيد من حجج الإدخال انظر
\citealt[pp.~vi, 24, 50--1]{Potter2004}.

ومع ذلك قلّ أن تسمح الرياضيات بالعناصر الأولية. ومن أسباب ذلك أن صوغ
نظرية المجموعات من غيرها أيسر، وإن كان الفرق \emph{يسيرًا جدًا}. وترد
أحيانًا حجة أعمق في استبعادها:
\begin{quote}
	وفقًا للاعتقاد بأن نظرية المجموعات هي أساس الرياضيات، ينبغي أن
	نستطيع احتواء الرياضيات كلها بمجرد الحديث عن المجموعات، ولذلك لا
	ينبغي لمتغيراتنا أن تتراوح على أشياء كالأبقار والخنازير.
	%لكن إذا كانت $C$ بقرة، فإن $\{C\}$ مجموعة، لكنها ليست شيئًا رياضيًا مشروعًا.
	\citep[p.~8]{Kunen1980}
\end{quote}
فالعناية بالتطبيق تدعو إلى \emph{إدخال} العناصر الأولية، والعناية بغاية
تأسيسية اختزالية، أي رد الرياضيات إلى نظرية المجموعات المحضة، تدعو إلى
\emph{استبعادها}. وشيء من حب الاختصار يدعو إلى الاستبعاد أيضًا.

وسنسلك الطريق الأقصر. وله أيضًا مسوغ تعليمي: فالمقصود تعريف القارئ
بأصول نظرية المجموعات التي يحتاج إليها لابتداء النظر في فلسفتها، وأكثر
البحوث الفلسفية في الباب تتناول نظريات صيغت \emph{من دون} عناصر أولية.
  فيكون الكتاب أنفع إذا قارب تلك البحوث فيما يعرضه.

\end{document}


